{\scriptsize\setlength{\tabcolsep}{3pt}
\begin{longtable}{rcccccp{5.7cm}}
\caption{Casos de teste de \texttt{stack\_reset.vwf}: reset assíncrono.}
\label{tab:caso-stack-reset}\\
\hline $t$ (ns) & \texttt{nrst} & \texttt{stack\_in} & \texttt{stack\_push} & \texttt{stack\_pop} & \texttt{stack\_out} & Observação \\ \hline \endfirsthead
\hline $t$ (ns) & \texttt{nrst} & \texttt{stack\_in} & \texttt{stack\_push} & \texttt{stack\_pop} & \texttt{stack\_out} & Observação \\ \hline \endhead
0--20 & 0 & 0 & 0 & 0 & 0 &  \\
20--40 & 1 & 11 & 1 & 0 & 11 & push 11 \\
40--60 & 1 & 22 & 1 & 0 & 22 & push 22 \\
60--80 & 1 & 33 & 1 & 0 & 33 & push 33 \\
80--85 & 1 & 33 & 0 & 0 & 33 & sem operação \\
85--100 & 0 & 33 & 0 & 0 & 0 & reset assíncrono, fora de qualquer borda \\
100--123 & 0 & 44 & 1 & 0 & 0 & ainda em reset; push sobre a borda: não empilha \\
123--140 & 1 & 44 & 1 & 0 & 44 & solta o reset \\
140--160 & 1 & 77 & 1 & 0 & 77 & push 77 (a borda já empilha) \\
160--180 & 1 & 77 & 0 & 1 & 44 & pop (o topo vira 0: nada sobrou do que veio antes do reset) \\
180--200 & 1 & 77 & 0 & 1 & 0 & pop \\
200--212 & 1 & 77 & 0 & 0 & 0 & sem operação \\
212--215 & 0 & 77 & 0 & 0 & 0 & pulso de reset de 3 ns, sem borda \\
215--220 & 1 & 77 & 0 & 0 & 0 & fim do pulso \\
220--240 & 1 & 5 & 1 & 0 & 5 & push 5 \\
240--260 & 1 & 5 & 0 & 0 & 5 & sem operação \\
\hline
\end{longtable}}
\vspace{-0.9cm}
\fonte{Autoria própria. Saídas do fim de cada intervalo, conforme o modelo de referência e a simulação do Quartus II.}
